\section{Introduction}
\label{sec:intro}

Retrieval-Augmented Generation (RAG) \cite{Base2} has become a prominent method for augmenting Large Language Models (LLMs) with external knowledge resources, enabling them to generate more accurate outputs and thereby expanding their practical utility \cite{Hallucination1,Hallucination2}. However, conventional RAG approaches, which primarily employ dense vector retrieval \cite{Base1,Base3,Base4} for identifying relevant text segments, often exhibit limitations when confronted with intricate tasks that necessitate multi-hop inference \cite{Benchmark2} or the capture of long-range dependencies \cite{Benchmark1}. Knowledge Graphs (KGs), characterized by their structured representation of interconnected entities and relationships, present a compelling alternative. Thus, research efforts are increasingly focused on developing KG-RAG systems, which aims to improve RAG performance using graph-based text representation \cite{KG4,KG7,KG9}.

Usually, a KG-RAG system operates by initially transforming source documents into graph-based representations\footnote{Unlike previous work leveraging established KGs \cite{KG1,KG2,LLM-Traverse7,KGQA1}, this study explores KG-RAG systems that generate their knowledge graph directly from the source text, like GraphRAG \cite{Main3} and LightRAG \cite{Main2}.}. This often involves leveraging LLMs to extract entities and relationships, which are then represented as nodes and edges within the graph. Information retrieval is subsequently performed by traversing this graph to identify nodes relevant to the query. To achieve this, traditional graph search algorithms, like PageRank \cite{PPR,Main3}, and deep learning approaches, like GNN \cite{GNN1}, could be employed for traversal and answer extraction \cite{Main5,KG3,Main9,Cot5}, these methods often fail to capture the nuanced semantic relationships embedded within the graph, leading to unsatisfactory system effectiveness \cite{Main6,KG4,LLM-Traverse5}. By contrast, the LLM-guided knowledge graph traversal approaches demonstrate notable advantages \cite{KG1,KG2,KG5}, but these methods lead to substantially increased LLM invocation costs and a significant rise in irrelevant noise information. Therefore, \textbf{the effective synergy between system effectiveness and cost efficiency presents a major challenge for the practical deployment and scalability of KG-RAG systems}.

To this end, we propose \ours—\textbf{Retrieve and Memorize}, a novel LLM-guided KG-RAG system with memorization. As depicted in \Cref{fig:overall-workflow}, \ours~ first constructs a KG from unstructured texts and subsequently employs two key innovative modules: (1) \textit{LLM-Guided Knowledge Graph Traversal with Node Exploration and Exploitation}. Beyond logical reasoning and semantic understanding capabilities of the LLM, \ours~expands traversal paths using node exploration and node exploitation. This approach enables the LLM to progressively perform single-node expansion and traversal to approximate answers, achieving high-precision retrieval with more hops without requiring large-scale searches like beam search-based \cite{Beam-Search} methods. (2) \textit{Memory Replay for Efficient Graph Traversal}. It memorizes traversal experience within KG edge embeddings, mirroring the way LLMs "memorize" world knowledge within their parameters, but in a train-free manner. It efficiently stores and recalls experience, eliminating the need to consult the LLM, thus substantially improving efficiency while preserving system effectiveness. We theoretically show that for semantically similar queries, this memorization enables edge embeddings to store sufficient traversal experience regarding these queries, facilitating effective memory replay.

\begin{figure}[h!]
    \centering
    \includegraphics[width=0.99\textwidth]{Resources/workflow.pdf}
    \caption{Overall Workflow. \ours~ constructs a KG from unstructured text. It memorizes LLM-guided traversal path, enabling fast retrieval when encountering similar queries subsequently.}
    \label{fig:overall-workflow}
\end{figure}

To evaluate \ours, we conduct extensive experiments across various benchmark datasets and LLM backbones. The experimental results demonstrate that \ours~exhibits a clear advantage over competing baseline approaches, achieving performance gains of 5\% to 10\% while simultaneously reducing the average cost per query by approximately 50\%. Our analysis reveals that the key strengths of \ours~ lie in its inherent self-correction capability and robust memory retention. It can leverage LLMs to self-correct erroneous search paths while maintaining stable memory capabilities even when handling large-scale updates. To sum up, our work makes the following core contributions.
\begin{itemize}[leftmargin = 15pt]
    \item We call attention to the challenge of achieving effective synergy between system effectiveness and cost efficiency, hindering the practical utility of KG-RAG systems.
    \item We introduce \ours, which employs LLM-guided KG traversal featuring node exploration, node exploitation, and, most notably, memory replay, to improve both system effectiveness and cost efficiency.
    \item We theoretically and experimentally confirm the effectiveness of our graph traversal with memory replay, demonstrating \ours's superiority over existing baselines across various benchmark datasets and LLM backbones.
\end{itemize}

